\input{../tex-inputs/latex-leseansicht-vorspann}

\section[Arthur Schnitzler und Felix Salten an Gustav Schwarzkopf, 29. 8. 1900]{L04518 Arthur Schnitzler und Felix Salten an Gustav Schwarzkopf, 29. 8. 1900}
\nopagebreak\mylabel{L04518v}
\rehead{ }\normalsize\beginnumbering\briefempfaengerindex{Schwarzkopf, Gustav@\textsc{Schwarzkopf, Gustav}!zzzSalten, Felix@\emph{von Felix Salten}!1900-08-291@{29. 8. 1900}|(be}\briefempfaengerindex{Schwarzkopf, Gustav@\textsc{Schwarzkopf, Gustav}!zzzSchnitzler, Arthur@\emph{von Arthur Schnitzler}!1900-08-291@{29. 8. 1900}|(be}
\toendnotes[C]{\smallbreak\pagebreak[2]}
\correspDesc{Versand  durch Arthur Schnitzler, Felix Salten am 29. 8. 1900 in Meran
\newline{}Übermittlung  am 29. 8. 1900 in Meran
\newline{}Zustellung  am 31. 8. 1900 in Wien
\newline{}Erhalt  durch Gustav Schwarzkopf im Zeitraum [31. 8. 1900 – 3. 9. 1900?] in Wien}\toendnotes[C]{\smallbreak}
\Standort{DLA, A:Schnitzler, HS.1985.1.1897.}
\physDesc{Bildpostkarte, 133 Zeichen
\newline{}Handschrift Arthur Schnitzler: Bleistift, deutsche Kurrent
\newline{}Handschrift Felix Salten: Bleistift, lateinische Kurrent
\newline{}Versand: 1) Stempel: »\nobreak{}\oindex{Meran@\textbf{Meran}|pwk}Meran, 29. 8. 00, 7\nobreak{}«.   2) Stempel: »\nobreak{}\oindex{Wien@\textbf{Wien}|pwk}Wien 1/1 1, 31. 8. 1900, 8–9 1/2 V., bestellt\nobreak{}«. }\toendnotes[C]{\smallbreak}
\pstart
           \noindent{}\centering{}{\pb}\textcolor{gray}{\textbf{Meran}}\oindex{Meran@\textbf{Meran}|pw}\hspace*{2.5em}\textcolor{gray}{\textbf{Gilfpromenade}}\oindex{Gilfpromenade@\textbf{Gilfpromenade}|pw}\pend
           \pstart{}{\pb}Herrn \textsc{Gustav Schwarzkopf}\pend{}\pstart{}\textsc{Wien}\oindex{Wien@\textbf{Wien}|pw}\pend{}\pstart{}\textsc{I. Tiefer Graben 23}\oindex{Wien@\textbf{Wien}!I., Innere Stadt@\textbf{I., Innere Stadt}!Tiefer Graben 23@\textbf{Tiefer Graben 23}|pw}\pend{}{\bigskip}\vspace{1em}
\pstart
           \noindent{}{\pb}Viele herzliche Grüße!\pend
           \pstart \spacefill\mbox{ArthSchn.}\pend{}\selectlanguage{ngerman}\vspace{1em}
\pstart
           \noindent{}{[}hs. Salten:{]} Herzl. Grüße\pend
           \pstart \spacefill\mbox{Salten}\pend{}\selectlanguage{ngerman}\vspace{1em}
\pstart
           {[}hs. Schnitzler:{]} 29.\,8.\,900\pend
           \vspace{0.5em}
\pstart
           (\label{K_L04518-1v}\edtext{P. G.\pwindex{Goldmann, Paul 31.\,1.\,1865 Breslau – 25.\,9.\,1935 Wien@\textsc{Goldmann, Paul} (31.\,1.\,1865 Breslau – 25.\,9.\,1935 Wien)|pw} blieb in Trafoi\oindex{Trafoi@\textbf{Trafoi}|pw} zurück}{\lemma{\textnormal{\emph{P. G. … zurück}}}\Cendnote{\textnormal{Siehe A. S.: \emph{Tagebuch}, 28. 8. 1900.}}}\label{K_L04518-1}.)\pend
           \selectlanguage{ngerman}\endnumbering\briefempfaengerindex{Schwarzkopf, Gustav@\textsc{Schwarzkopf, Gustav}!zzzSalten, Felix@\emph{von Felix Salten}!1900-08-291@{29. 8. 1900}|)be}\briefempfaengerindex{Schwarzkopf, Gustav@\textsc{Schwarzkopf, Gustav}!zzzSchnitzler, Arthur@\emph{von Arthur Schnitzler}!1900-08-291@{29. 8. 1900}|)be}\mylabel{L04518h}
\begin{anhang}
\end{anhang}\newcommand{\dateiname}{L04518}\newcommand{\titel}{Arthur Schnitzler und Felix Salten an Gustav Schwarzkopf, 29. 8. 1900}\newcommand{\editorInnen}{Herausgegeben von Jahnke, SelmaMüller, Martin Anton}\input{../tex-inputs/latex-leseansicht-abspann}
